\documentclass{stex}

\libusepackage{naproche}
\libusepackage{copyright}
\libinput{preamble}

\begin{document}

\usemodule[libraries/meta]{preliminaries.ftl}


\section{Notions}\label{sec:notions}

\subsection{Objects}\label{sec:objects}

\emph{Objects} are mathematical entities that are ``small'' enough to be
contained in classes.

\begin{fakeforthel}
  \begin{signature*}
    A \emph{mathematical object} is a notion.
    Let an \emph{object} stand for a mathematical object.
    Let an \emph{element} stand for a mathematical object.
  \end{signature*}
\end{fakeforthel}

\subsection{Ordered Pairs}\label{sec:ordered-pairs}

\begin{fakeforthel}
  \begin{signature*}
    Let $a,b$ be objects.
    $\emph{(a,b)}$ is an object.
    Let the \emph{ordered pair of $a$ and $b$} stand for $(a,b)$.
  \end{signature*}
\end{fakeforthel}

\begin{fakeforthel}
  \begin{axiom*}[forthel,title=Ordered Pair Extensionality]
    Let $a,a',b,b'$ be objects.
    If $(a, b) \Eq (a', b')$ then $a \Eq a'$ and $b \Eq b'$.
  \end{axiom*}
\end{fakeforthel}


\subsection{Classes and Sets}\label{sec:classes-and-sets}

\begin{fakeforthel}
  \begin{signature*}
    A \emph{class} is a notion.
    Let a \emph{collection} stand for a class.
  \end{signature*}
\end{fakeforthel}

\begin{fakeforthel}
  \begin{axiom*}[title=Class Extensionality]
    Let $A,B$ be classes.
    If every element of $A$ is an element of $B$ and every element of $B$ is an element of $A$ then $A \Eq B$.
  \end{axiom*}
\end{fakeforthel}

Sets are considered to be classes that are ``small'' enough to be contained in
other classes:

\begin{fakeforthel}
  \begin{definition*}
    A \emph{set} is a class that is an object.
  \end{definition*}
\end{fakeforthel}

\begin{sparagraph}[style=symdoc,for={Elem,NotElem}]
  \begin{fakeforthel}
    \begin{signature*}
      Let $A$ be a class.
      A \emph{element of $A$} is an object.
      Let a \emph{member of $X$} stand for an element of $X$.
      Let $x$ \emph{belongs to $X$} stand for $x$ is an element of $X$.
      Let $X$ \emph{contains $x$} stand for $x$ is an element of $X$.
      Let $x$ is \emph{contained in $X$} stand for $x$ is an element of $X$.
      Let $x$ \emph{lies in $X$} stand for $x$ is an element of $X$.
      Let $x$ is \emph{in $X$} stand for $x$ is an element of $X$.
      Let $\emph{a \Elem A}$ stand for $a$ is an element of $A$.
      Let $\emph{a \NotElem A}$ stand for $\NOT a \Elem A$.
    \end{signature*}
  \end{fakeforthel}
\end{sparagraph}


\subsection{Comprehension Terms}\label{sec:comprehension-terms}

In \ForTheL we have three types of comprehension terms:
\begin{itemize}
  \item ``Enumeration classes'', i.e. terms of the form
  $\EClass{a_1,\dots,a_n}$, that represent the class that contains precisely the
  objects $a_1,\dots,a_n$.
  Typical examples are $\EClass{4,7,11}$ (i.e. the class consisting of the
  numbers $4$, $7$ and $11$) or $\EClass{x,y+1,z-x}$ (i.e. the class consisting
  of the elements $x$, $y+1$ and $z-x$ for some fixed objects $x$, $y$ and $z$).
  \item ``Comprehension classes'', i.e. terms of the form
  $\CClass{\tau}{\varphi}$, that represent the class of all objects of the form
  $\tau$ such that the components of $\tau$ satisfy the condition $\varphi$.
  Typical examples are $\CClass{p}{p \text{ is an odd prime number}}$ (i.e the
  class of all odd prime numbers) or $\CClass{(p,n)}{n \in \mathbb{N}
  \text{ and } p \text{ is a prime divisor of } n}$ (i.e. the class of all pairs
  $(p,n)$ of natural numbers such that $p$ is a prime divisor of $n$).
  \item ``Separation classes'', i.e. terms of the form
  $\SClass{\tau}{A}{\varphi}$, that represent the subclass of a class $A$ of all
  objects of the form $\tau$ such that the components of $\tau$ satisfy the
  condition $\varphi$.
  Typical examples are $\SClass{p}{\mathbb{P}}{p \text{ is odd}}$ (i.e the
  class of all odd prime numbers) or $\SClass{(p,n)}{\mathbb{N}\times\mathbb{N}}
  {p \text{ is a prime divisor of } n}$ (i.e. the class of all pairs
  $(p,n)$ of natural numbers such that $p$ is a prime divisor of $n$).
\end{itemize}

\begin{sparagraph}[style=symdoc,for=EClass]
  \begin{signature*}
    Let $a_1, \dots, a_n$ be objects.
    $\fakeemph{\EClass{a1, \dots, a_n}}$ is a class.
  \end{signature*}

  \begin{axiom*}
    Let $a_1, \dots, a_n$ be objects.
    Then $a \Elem \EClass{a_1, \dots, a_n}$ iff $a \Eq a_1$ or \dots{} or $a \Eq a_n$.
  \end{axiom*}
\end{sparagraph}

\begin{sparagraph}[style=symdoc,for=CClass]
  \begin{signature*}
    Let $\tau$ be a term and $\varphi$ be a formula.
    $\fakeemph{\CClass{\tau}{\varphi}}$ is a class.
  \end{signature*}

  \begin{axiom*}
    Let $\tau[x_1, \dots, x_n]$ be a term and $\varphi[x_1, \dots, x_n]$ be a formula.
    Then $a \Elem \CClass{\tau}{\varphi}$ iff
    $a$ is an object and there exist $a_1, \dots, a_n$ such that $\varphi[a_1, \dots, a_n]$ and $a \Eq \tau[a_1, \dots, a_n]$.
  \end{axiom*}

  \begin{proposition*}
    Let $\varphi[x]$ be a formula.
    Then $a \Elem \CClass{x}{\varphi}$ iff $a$ is an object and $\varphi[a]$.
  \end{proposition*}
\end{sparagraph}

\begin{sparagraph}[style=symdoc,for=SClass]
  \begin{signature*}
    Let $\tau$ be a term, $A$ be a class and $\varphi$ be a formula.
    $\fakeemph{\SClass{\tau}{A}{\varphi}}$ is a class.
  \end{signature*}

  \begin{axiom*}
    Let $\tau[x_1, \dots, x_n]$ be a term, $A$ be a class and $\varphi[x_1,\dots,x_n]$ be a formula.
    Then $a \Elem \SClass{\tau}{A}{\varphi}$ iff
    $a \Elem A$ and there exist $a_1, \dots, a_n$ such that $\varphi[a_1, \dots, a_n]$ and $a \Eq \tau[a_1, \dots, a_n]$.
  \end{axiom*}

  \begin{proposition*}
    Let $A$ be a class and $\varphi[x]$ be a formula.
    Then $a \Elem \SClass{x}{A}{\varphi}$ iff
    $a \Elem A$ and $\varphi[a]$.
  \end{proposition*}

  \begin{axiom*}
    Let $\tau$ be a term, $A$ be a class and $\varphi$ be a formula.
    If $A$ is a set then $\SClass{\tau}{A}{\varphi}$ is a set.
  \end{axiom*}
\end{sparagraph}


\subsection{Maps and Functions}\label{sec:maps-and-functions}

\begin{fakeforthel}
  \begin{signature*}
    A \emph{map} is a notion.
  \end{signature*}
\end{fakeforthel}

\begin{fakeforthel}
  \begin{axiom*}[title=Map Extensionality]
    Let $f,g$ be maps.
    If $\Dom(f) \Eq \Dom(g)$ and $f(a) \Eq g(a)$ for all $a \Elem \Dom(f)$ then $g \Eq g$.
  \end{axiom*}
\end{fakeforthel}

Similar to the notion of sets which are considered to be ``small'' classes, we
consider \emph{functions} to be ``small'' maps.

\begin{fakeforthel}
  \begin{definition*}
    A \emph{function} is a map that is an object.
  \end{definition*}
\end{fakeforthel}

\begin{forthel}
  \begin{axiom*}
    Let $f$ be a map.
    $f$ is a function iff $\Dom(f)$ is a set.
  \end{axiom*}
\end{forthel}

\begin{sparagraph}[style=symdoc,for=Dom]
  \begin{fakeforthel}
    \begin{signature*}
      Let $f$ be a map.
      $\emph{\Dom(f)}$ is a class.
      Let the \emph{domain of $f$} stand for $\Dom(f)$.
      Let $f$ is \emph{defined on $X$} stand for $\Dom(f) \Eq X$.
    \end{signature*}
  \end{fakeforthel}
\end{sparagraph}

\begin{fakeforthel}
  \begin{signature*}
    Let $f$ be a map and $a \Elem \Dom(f)$.
    $\emph{f(a)}$ is an object.
    Let the \emph{value of $f$ at $x$} stand for $f(x)$.
    Let $\emph{f(x,y)}$ stand for $f((x,y))$.
  \end{signature*}
\end{fakeforthel}


\subsection{$\lambda$-Terms}\label{sec:lambda-terms}

In low-level map definitions, \ForTheL allows to define maps via
$\lambda$-terms:

\begin{sparagraph}[style=symdoc,for=Map]
  \begin{signature*}
    Let $x$ be a variable, $A$ be a class and $\tau$ be a term.
    $\fakeemph{\Map{x}{A}\tau}$ is a map.
  \end{signature*}

  \begin{axiom*}
    Let $x$ be a variable, $A$ be a class and $\tau[x]$ be a term.
    Then $\Dom(\Map{x}{A}\tau) \Eq A$.
  \end{axiom*}

  \begin{axiom*}
    Let $x$ be a variable, $A$ be a class and $\tau[x]$ be a term and $a \Elem A$.
    Then $(\Map{x}{A}\tau)(a) \Eq \tau[a]$.
  \end{axiom*}
\end{sparagraph}


\ifinputref\else\printlicense[CcByNcSa]{Marcel Schütz (2026)}\fi

\end{document}
